\par\Needspace{8\baselineskip}\bigskip
\pdfbookmark[3]{Darstellung, 13}{pallas-II-article13}
\begin{center}13.\end{center}
\noindent\hspace*{\parindent}Nach der üblichen Konvention ist die Epoche der mittleren Länge für die Zeit \(0\) gleich \(L-nt\). Nach dieser Konvention ergäbe sich also
\[
\begin{aligned}
&\mathrm{d}(\text{Epoche der mittl. Länge})=-t\mathrm{d}n+(1-\cos i)\mathrm{d}\AscendingNode\\
&\quad-(2ar+aa\cos\varphi\mathbin{.}\operatorname{tang}\tfrac12\varphi\mathbin{.}\cos(v-\tilde\omega))Tn\mathrm{d}t\\
&\quad+\dfrac{ar\operatorname{tang}\tfrac12\varphi\mathbin{.}\sin(v-\tilde\omega)\mathbin{.}(2+e\cos(v-\tilde\omega))}{\cos\varphi}Vn\mathrm{d}t,
\end{aligned}
\]
wobei man für\raisebox{-.45ex}{-} \(\mathrm{d}n\) den oben gefundenen Wert einsetzen könnte. So enthielte die momentane Änderung der Epoche jedoch einen Term, der wegen des Faktors \(t\) unbegrenzt wachsen könnte. Um diese Schwierigkeit zu vermeiden, definieren wir die Epoche der mittleren Länge durch \(L-\int n\mathrm{d}t\). Dies stimmt, wie man sieht, bei der elliptischen Bewegung mit der üblichen Konvention überein. Bezeichnen wir also die Epoche der mittleren Länge für die Zeit null mit \(\epsilon\), so erhalten wir
\[
\begin{array}{rl}
70)\quad&\begin{aligned}[t]\mathrm{d}\epsilon&=(1-\cos i)\mathrm{d}\AscendingNode\\
&\quad-(2ar+aa\cos\varphi\mathbin{.}\operatorname{tang}\tfrac12\varphi\mathbin{.}\cos(v-\tilde\omega))Tn\mathrm{d}t\\
&\quad+\dfrac{ar\operatorname{tang}\tfrac12\varphi\mathbin{.}\sin(v-\tilde\omega)\mathbin{.}(2+e\cos(v-\tilde\omega))}{\cos\varphi}Vn\mathrm{d}t,\end{aligned}
\end{array}
\]
und die mittlere Länge für die Zeit \(t\) ergibt sich aus der Formel
\[
\begin{array}{rl}
71)\quad&L=\epsilon+\int n\mathrm{d}t.
\end{array}
\]
